\documentclass[11pt]{article}
\usepackage[margin=1in]{geometry}
\usepackage{fontspec}
\usepackage{titlesec}
\usepackage{color}
\usepackage{hyperref}
\usepackage{enumitem}
\usepackage{xcolor}
\usepackage{setspace}
\usepackage[absolute,overlay]{textpos}


% Fonts and layout
\setmainfont{Times New Roman}
\titleformat{\section}{\large\bfseries\color{brown}}{}{0em}{}[\titlerule]
\titleformat{\subsection}{\normalsize\bfseries\color{Sepia}}{}{0em}{}

% Colors
\definecolor{Sepia}{RGB}{112, 66, 20}
\definecolor{Brown}{RGB}{92, 64, 51}

\hypersetup{
    colorlinks=true,
    linkcolor=Brown,
    urlcolor=Sepia
}

\pagestyle{empty}
\setlength{\parskip}{0.5em}
\setlength{\parindent}{0pt}

\begin{document}

\begin{textblock*}{5cm}(15cm,1cm)
  \raggedleft
  \textit{Last updated: \today}
\end{textblock*}

\begin{center}
    {\LARGE \textbf{Curriculum Vitae}}\\
    \vspace{0.5em}
    \textbf{Kanta Ogawa}\\
    Email: \href{mailto:ko2406@nyu.edu}{ko2406@nyu.edu}
\end{center}

\vspace{1em}

\section*{Education}
\begin{itemize}[leftmargin=1.5em]
  \item \textbf{Ph.D. in Economics}, New York University, USA \\
  2025 -- 2031(expected)
  \item \textbf{M.A. in Economics}, University of Tokyo, Japan \\
  2023 -- 2025
  \item \textbf{B.A. in Economics}, University of Tokyo, Japan \\
  2019 -- 2023
\end{itemize}

\section*{Research Interests}
Macro-labor, heterogeneous household behavior, taxation, and social security systems

\section*{Working Paper}
\textbf{\href{https://arxiv.org/pdf/2503.17917}{Part-Time Penalties and Heterogeneous Retirement Decisions}}

Revise and resubmit at \textit{Macroeconomic Dynamics}.

\textit{Abstract:} Older male workers exhibit substantial heterogeneity in retirement behavior across occupations and respond differently to retirement policy reforms. A key source of this heterogeneity is the part-time penalty—the reduction in hourly wages when workers switch from full-time to part-time employment. Workers in occupations with higher part-time penalties are more likely to retire earlier, as observed in data from Japan and the United States. To explain this pattern, I develop a general equilibrium model, demonstrating that greater asset accumulation and more generous pension benefits raise reservation wages and contribute to this earlier retirement. Using the Japanese Panel Study of Employment Dynamics, I show that standard retirement policy reforms—eliminating the earnings test, extending the pension eligibility age, and reducing pension benefits—have modest or negative effects on aggregate effective labor and output. In contrast, expanding income tax credits and exempting pension benefits from income taxation increase aggregate effective labor by about 1\% across occupations, raise output by 1.197\% and 0.5492\%, respectively, and improve long-term welfare.
\section*{Presentations}

\textbf{2025}
\begin{itemize}[leftmargin=1.5em]
   \item Society for Computational Economics (Santiago, Chile)
   \item RCEA International Conference on Economics, Econometrics, and Finance (Jersey City, USA)
   \item Spring Meeting of the Japanese Economic Association (Nagoya, Japan)
\end{itemize}

\section*{Research Experience}
\begin{itemize}[leftmargin=1.5em]
  \item \textbf{May 2024 -- December 2024}, National Graduate Institute for Policy Studies \\
  Research Assistant to Prof. Sagiri Kitao and Prof. Daiji Kawaguchi
  \item \textbf{February 2024 -- May 2024}, University of Tokyo \\
  Research Assistant to Prof. Sagiri Kitao
  \item \textbf{August 2023 -- June 2025}, Institute of Social Science, University of Tokyo \\
  Research Assistant to Prof. Ayako Kondo and Prof. Daiji Kawaguchi
  \item \textbf{February 2023 -- March 2023}, University of Tokyo \\
  Research Assistant to Prof. Sagiri Kitao
  \item \textbf{August 2022 -- December 2026}, UTokyo Economic Consulting Inc. \\
  Research Assistant
\end{itemize}

\section*{Teaching Experience}
\begin{itemize}[leftmargin=1.5em]
  \item \textbf{October 2023 -- January 2024}, University of Tokyo \\
  Teaching Assistant for \textit{Macroeconomics II} (core Ph.D. course) taught by Prof. Sagiri Kitao
\end{itemize}

\section*{Fellowships}
\begin{itemize}[leftmargin=1.5em]
  \item \textbf{2025 -- Present}: MacCracken Fellowship, New York University
  \item \textbf{2025 -- Present}: Heiwa-Nakajima Foundation
  \item \textbf{2024 -- 2025}: The World-leading Innovative Graduate Study: Economics
  \item \textbf{2023 -- 2025}: The World-leading Innovative Graduate Study: Advanced Business Law Program
\end{itemize}

\section*{Programming Skills}
Python, R, MATLAB, Fortran, Stata

\section*{Languages}
\begin{itemize}[leftmargin=1.5em]
  \item Japanese (native)
  \item English (fluent)
\end{itemize}

\section*{Citizenship}
Japan



\end{document}
